% Compile: lualatex demo && bibtex demo && lualatex demo && lualatex demo
\documentclass[10pt]{article}
\usepackage[journal]{cogra}   % try [twocol], [thesis], or add aged
\cograrunning{E. G. Pellham and D. R. Okonkwo}{Goal-Box Illumination and Running Speed}

\begin{document}
\cogratitle{Running Speed in a Straight Alley as a Function of Goal-Box Illumination}
\cograauthor{Edwin G. Pellham and Dorothy R. Okonkwo}
\cograaffil{Department of Psychology, Halvorsen College}
\cogranote{This article is fictional. It was written to show the cogra
style, and its authors, institution, data and references do not exist.}
\cogramaketitle

\begin{cograsummary}
Two groups of hungry rats were run to food in a straight alley whose goal
box was either dimly or brightly lit. Running speed rose over trials in
both groups but reached a lower level when the goal box was bright. The
result is discussed in terms of a competing tendency to avoid the lighted
box.
\end{cograsummary}

It is commonly supposed that the speed with which an animal approaches a
goal depends on what is found there, and that the surroundings of the goal
matter little once the animal has been fed in them a few times. The
supposition has not often been put to test. \citet{brandt1949} reported
that rats ran more slowly to a goal box painted white than to one painted
black, but attributed the difference to the reflectance of the box rather
than to its illumination, and the matter has been left there since
\citep{halloran1953}. The present experiment varied the illumination of
the goal box directly, holding its surfaces constant, to see whether the
lighting of the place in which food is found alters the speed of running
to it.\footnote{A symbol, not a number, marks each note, as was common in
journals of the period.}

\section{Method}
\subsection{Subjects}
The subjects were 24 male albino rats, about 100 days old at the start of
the experiment, maintained at 85\% of their free-feeding weight
throughout. They were assigned at random to two groups of 12.

\subsubsection{Apparatus and procedure}
The alley was a straight runway 6 ft long, painted a uniform flat gray,
with a start box at one end and a goal box at the other. A photocell at
each end started and stopped an electric clock. For the dim group the goal
box was lit by a 7.5-w lamp behind a diffusing screen; for the bright
group the same lamp was replaced by one of 60 w. Each rat received 10
trials a day for 6 days, with a pellet of food in the goal box on every
trial. Running time was converted to speed in feet per second.

\section{Results}
Mean running speed in each block of 10 trials is shown in Fig.~1, and the
values for the final block are given in Table~1. Both groups ran faster as
training went on, but the bright group reached a lower level than the dim
group from the third block onward \citep[cf.][]{brandt1949}.

\begin{figure}[t]
\centering
\begin{tikzpicture}
\begin{axis}[cogra, width=0.82\linewidth, height=0.55\linewidth,
  xlabel={Blocks of 10 trials}, ylabel={Running speed (ft/sec)},
  xmin=0.5, xmax=6.5, ymin=0, ymax=3.2, xtick={1,...,6},
  legend pos=north west]
\addplot coordinates {(1,0.62)(2,1.18)(3,1.71)(4,2.14)(5,2.41)(6,2.56)};
\addplot coordinates {(1,0.58)(2,1.02)(3,1.33)(4,1.55)(5,1.68)(6,1.74)};
\legend{Dim goal box,Bright goal box}
\end{axis}
\end{tikzpicture}
\caption{Mean running speed as a function of blocks of trials for the
dim and bright groups, plotted in the open and filled circle convention
of the period.}
\end{figure}

\begin{table}[t]
\caption{Running speed and latency on the final block of trials}
\centering\small
\begin{tabular}{lcc}
\toprule
Group & Speed (ft/sec) & Start latency (sec) \\
\midrule
Dim goal box    & 2.56 & 1.4 \\
Bright goal box & 1.74 & 2.1 \\
\bottomrule
\end{tabular}
\end{table}

\section{Discussion}
The lighting of the goal box, which bore no relation to the food found in
it, lowered the speed with which the animals ran to it. One account
supposes that a bright place arouses in the rat a tendency to withdraw,
which competes with the approach set up by feeding
\citep{marchetti1955}; another, that bright light reduces the value of the
reward itself \citep{strand1957}. The two can be separated by varying the
illumination of the alley while holding that of the goal box constant, and
an experiment of this kind is now in progress.

\bibliographystyle{apalike}
\bibliography{demo}
\end{document}
